% ============================================================
%  slide_ch1.tex — Chương I: Mở đầu
%  5 frames
% ============================================================
\section{Chương I: Mở đầu}

% ------ Frame 1: Bối cảnh ------
\begin{frame}{Bối cảnh: Giao thông đô thị Việt Nam}
  \begin{columns}[T]
    \column{0.56\textwidth}
    \begin{itemize}
      \item Ùn tắc giao thông gây thiệt hại \alert{hàng tỷ đồng/ngày} tại Hà Nội và TP.HCM
      \vspace{4pt}
      \item Hệ thống đèn hiện tại: \textbf{Fixed-Time Control}
        \begin{itemize}
          \item Chu kỳ đèn cố định, không thay đổi
          \item Không phản ứng theo lưu lượng thực tế
          \item Lãng phí thời gian xanh khi đường vắng
        \end{itemize}
      \vspace{4pt}
      \item Đặc thù riêng của VN: xe máy chiếm \alert{60--70\%} tổng lưu lượng
    \end{itemize}

    \column{0.40\textwidth}
    \begin{block}{Phân bố phương tiện VN}
      \begin{center}
      \begin{tabular}{lc}
        \toprule
        \textbf{Loại xe} & \textbf{Tỉ lệ}\\
        \midrule
        Xe máy     & $\sim$60\%\\
        Ô tô       & $\sim$30\%\\
        Xe buýt    & $\sim$5\%\\
        Xe tải     & $\sim$5\%\\
        \bottomrule
      \end{tabular}
      \end{center}
    \end{block}
    \vspace{6pt}
    \begin{alertblock}{}
      \centering Chưa có nghiên cứu DRL\\ phù hợp đặc thù VN!
    \end{alertblock}
  \end{columns}
\end{frame}

% ------ Frame 2: Vấn đề & Giải pháp ------
\begin{frame}{Vấn đề và Hướng tiếp cận}
  \begin{columns}[T]
    \column{0.48\textwidth}
    \begin{alertblock}{Vấn đề với Fixed-Time}
      \begin{itemize}
        \item Không nhận biết lưu lượng thực tế
        \item Chu kỳ cứng nhắc, không thích ứng
        \item Thiết kế theo kinh nghiệm, không tối ưu
        \item Nghiên cứu DRL hiện có chủ yếu cho giao thông Châu Âu / Bắc Mỹ
      \end{itemize}
    \end{alertblock}

    \column{0.48\textwidth}
    \begin{exampleblock}{Giải pháp: Deep RL (DQN)}
      \begin{itemize}
        \item Agent \textbf{tự học} từ dữ liệu mô phỏng
        \item Thích ứng pha đèn theo thời gian thực
        \item Không cần biết mô hình giao thông tường minh
        \item Tích hợp đặc thù VN (PCU xe máy, phân bố phương tiện)
      \end{itemize}
    \end{exampleblock}
  \end{columns}

  \vspace{10pt}
  \begin{center}
    \begin{tikzpicture}
      \node[draw, fill=vnred!20, rounded corners, minimum width=3cm, minimum height=0.55cm] (ft)
        at (0,0) {\small Fixed-Time: chu kỳ cố định};
      \node[draw, fill=techgreen!25, rounded corners, minimum width=3cm, minimum height=0.55cm] (dqn)
        at (6,0) {\small DQN Agent: thích ứng thực tế};
      \draw[->, thick, vnblue, line width=1.5pt] (ft) -- node[above]{\small\textbf{Thay thế bằng}} (dqn);
    \end{tikzpicture}
  \end{center}
\end{frame}

% ------ Frame 3: Mục tiêu ------
\begin{frame}{Mục tiêu dự án}
  \begin{enumerate}
    \setlength{\itemsep}{8pt}
    \item \textbf{Xây dựng môi trường mô phỏng} ngã tư Hà Nội bằng \textbf{SUMO}, tích hợp phân bố xe và hệ số PCU Việt Nam

    \item \textbf{Mô hình hóa bài toán} điều khiển đèn dưới dạng\\
          \textbf{Mô hình Quyết định Markov (MDP)} --- xác định rõ State, Action, Reward

    \item \textbf{Triển khai Double DQN + Dueling Network} để huấn luyện agent điều khiển đèn tự thích ứng theo lưu lượng thực tế

    \item \textbf{Đánh giá định lượng} DQN so với Fixed-Time 160s qua 5 chỉ số:\\
          hàng đợi, thời gian chờ, tốc độ, throughput, tổng reward
  \end{enumerate}
\end{frame}

% ------ Frame 4: Phạm vi & Cấu trúc ------
\begin{frame}{Phạm vi nghiên cứu và Cấu trúc báo cáo}
  \begin{columns}[T]
    \column{0.50\textwidth}
    \begin{block}{Phạm vi nghiên cứu}
      \begin{itemize}
        \item \textbf{Đối tượng}: Ngã tư 4 hướng, đô thị VN
        \item \textbf{Công cụ}: SUMO + TraCI API
        \item \textbf{Ngôn ngữ}: Python 3.9+ / PyTorch 2.2+
        \item \textbf{Kịch bản}: Hà Nội mẫu\\
              lưu lượng EW:NS = 2:1
        \item \textbf{Thời gian}: 1 giờ / episode
        \item \textbf{Huấn luyện}: 200.000 bước
      \end{itemize}
    \end{block}

    \column{0.46\textwidth}
    \begin{block}{Cấu trúc báo cáo}
      \begin{tabular}{ll}
        \textbf{Ch. I}   & Mở đầu \\[2pt]
        \textbf{Ch. II}  & Cơ sở lý thuyết \\[2pt]
        \textbf{Ch. III} & Phương pháp đề xuất \\[2pt]
        \textbf{Ch. IV}  & Thực nghiệm \& Đánh giá \\[2pt]
        \textbf{Ch. V}   & Kết luận \\
      \end{tabular}
    \end{block}

    \vspace{8pt}
    \begin{exampleblock}{Kết quả nổi bật}
      \centering
      $\downarrow$\textbf{51.6\%} hàng đợi\\
      $\downarrow$\textbf{74.8\%} thời gian chờ\\
      $\uparrow$\textbf{16.0\%} tốc độ lưu thông
    \end{exampleblock}
  \end{columns}
\end{frame}
